%! Least-Squares Regression — Unit 2, Lesson 2.4 — Guided Notes & Practice
% THE in-class document, in the MAIN IDEAS / NOTES shape (LESSON_SHAPE.md §1), at the COMPACT
% budget: the vocabulary box, then ONE guidednotes table of four rows, 14 problems, 2 pages.
%   I DO   : the teacher fills the definition lines and works the FIRST problem of each grid
%            (1, 4, 7, 11).
%   WE DO  : the class works the rest of each grid, students holding the pen; the last row,
%            Guided Practice, is worked entirely together in a fresh context.
%   YOU DO : nothing on this page — ../homework/ IS the individual practice.
% THE TRAP is problem 5 (r from R-Sq: the square root loses the sign). THE CRUX is problem 10,
% in the last instruction row: "their reaction time goes down 24 ms" — the slope read as a
% promise about every individual instead of a change in the PREDICTED value (EK DAT-1.H.2).
%
% THE DATA (instruction rows): 20 students, sleep x (hr, 4 to 10) vs reaction time y (ms).
%   xbar 7.2, s_x 1.5, ybar 310, s_y 45, r -0.80  ->  b = -24, a = 482.8, r^2 = 0.64.
%   Output consistent with n = 20: S = 27.74, SE(b) = 4.243, T = -5.66; SE(a) = 31.17, T = 15.49.
% GUIDED PRACTICE: 30 used Civics, age (yr, 0 to 10) vs price ($1000s): a = 24.60, b = -1.850,
%   R-Sq 81.0% -> r = -0.90. Here the intercept IS meaningful (a new car) — the contrast.
%
% PAGE LOCKSTEP with notes_key/: \blank -> \ans, \termblank -> \termans, the fourth argument of
% each \pcell. Nothing else differs. `\\` inside a cell ends the row: use \par.
\documentclass[12pt]{article}
\usepackage{apstats-article}
\usepackage{apstats-boxes}

\begin{document}

\pageheader{Unit 2, Lesson 2.4}{Guided Notes \& Practice}
% NAMESTRIP: no \namedateperiod here — the name row lives on the cover only.

\vspace{2pt}

% ── VOCABULARY (I do — students fill each term AS IT IS NAMED, never in advance) ──
\boxguard[16]
\begin{vocabbox}
\small
Fill in each term \textbf{as we name it} in the notes below.
\par
\termblank{Least-squares regression line}
\par
\termblank{Slope of the LSRL}
\par
\termblank{Coefficient of determination ($r^2$)}
\par
\end{vocabbox}

\small
\begin{guidednotes}
% ── ROW 1: the least-squares line from summary statistics ────────────────────
\mainidea[The best line:]{Least squares} &
The \textbf{least-squares regression line (LSRL)} is the one line that makes the sum of the
\blank{1.8cm} residuals as small as possible. It always passes through \blank{1.5cm}.

\vspace{2pt}
Twenty students: \textbf{sleep} $x$ (hours, 4 to 10) and \textbf{reaction time} $y$ (ms).
\ $\bar x = 7.2$, \ $s_x = 1.5$, \ $\bar y = 310$, \ $s_y = 45$, \ $r = -0.80$.
\par\vspace{2pt}
\stepnum{1}\ Slope: $b = r \cdot s_y / s_x$ --- always the same \blank{1.2cm} as $r$.\par
\stepnum{2}\ Intercept: $a = \bar y - b\,\bar x$, because the line goes through
$(\bar x, \bar y)$.\par
\stepnum{3}\ Write $\hat y = a + bx$ with the \blank{2.8cm} in place of $x$ and $y$.
\notesprompt{Build the sleep line.}
\begin{probgrid}{|Y|Y|Y|}
\pcell{1}{The slope $b$}{1.3cm}{} &
\pcell{2}{The intercept $a$}{1.3cm}{} &
\pcell{3}{The equation, with the variable names}{1.3cm}{} \\ \hline
\end{probgrid}
\\ \hline
% ── ROW 2: computer output and r^2 (split into two sub-rows at the prompt) ───
\mainidea[Computer output \&]{R-squared} &
{\centering\footnotesize\ttfamily
\begin{tabular}{@{} l r r r r @{}}
Predictor & Coef & SE Coef & T & P \\ \hline
Constant & 482.8 & 31.17 & 15.49 & 0.000 \\
Sleep & $-$24.00 & 4.243 & $-$5.66 & 0.000 \\
\end{tabular}\par
S = 27.74 \quad R-Sq = 64.0\% \quad R-Sq(adj) = 62.0\%\par}
\vspace{2pt}
Software, same 20 students: $a$ is the Coef in the \blank{2.0cm} row, $b$ the Coef in the
\blank{1.4cm} row. The \textbf{coefficient of determination} $r^2$ is the \blank{2.0cm} of the
\blank{1.8cm} in $y$ explained by the linear model with $x$.
\\
& \notesprompt{Read the output.}
\begin{probgrid}{|>{\hsize=0.75\hsize\linewidth=\hsize}X|>{\hsize=0.75\hsize\linewidth=\hsize}X|>{\hsize=1.5\hsize\linewidth=\hsize}X|}
\pcell{4}{Find $a$ and $b$. Do they match 1 and 2?}{1.4cm}{} &
\pcell{5}{R-Sq is 64.0\%. What is $r$?}{1.4cm}{} &
\pcell{6}{Interpret $r^2$ in context.}{1.4cm}{} \\ \hline
\end{probgrid}
\\ \hline
% ── ROW 3: interpreting slope and intercept — THE CRUX (two sub-rows) ────────
\mainidea[In context:]{Slope \& intercept} &
The \textbf{slope} $b$ is how much the \blank{1.9cm} $y$ changes for each \blank{1.6cm}
increase in $x$ --- a fact about the line, not a promise about every individual. The
\textbf{intercept} $a$ is the predicted $y$ when $x =$ \blank{0.7cm}; it means something only
if $x = 0$ makes sense \emph{and} is \blank{2.4cm}.
\par\vspace{3pt}
\fcolorbox{navy}{sky}{\parbox{\dimexpr\linewidth-2\fboxsep-2\fboxrule\relax}{\centering
\textbf{For each additional [unit of $x$], the \emph{predicted} [$y$] increases (decreases)
by [$|b|$ units].}}}
\\
& \notesprompt{Interpret the sleep line, $\widehat{\text{time}} = 482.8 - 24(\text{sleep})$.}
\begin{probgrid}{|Y|Y|}
\pcell{7}{Interpret the slope in context.}{1.7cm}{} &
\pcell{8}{Interpret the intercept. Is it meaningful here?}{1.7cm}{} \\ \hline
\pcell{9}{Jon slept 2 hours more than Ana. How do their \emph{predicted} times
compare?}{1.7cm}{} &
\pcell{10}{Maya writes: ``When a student sleeps one more hour, their reaction time goes down
24~ms.'' What is wrong?}{1.7cm}{} \\ \hline
\end{probgrid}
\\ \hline
% ── LAST ROW: GUIDED PRACTICE — we do, together, a fresh context ─────────────
\mainidea[Guided practice]{Used cars} &
A dealer recorded the \textbf{age} (years, from 0 to 10) and the \textbf{price} (thousands of
dollars) of 30 used Honda Civics and regressed price on age:
\par\vspace{2pt}
{\centering\footnotesize\ttfamily
\begin{tabular}{@{} l r r r r @{}}
Predictor & Coef & SE Coef & T & P \\ \hline
Constant & 24.60 & 0.950 & 25.89 & 0.000 \\
Age & $-$1.850 & 0.169 & $-$10.93 & 0.000 \\
\end{tabular}\par
S = 2.10 \quad R-Sq = 81.0\% \quad R-Sq(adj) = 80.3\%\par}
\\
& \notesprompt{We work these together.}
\begin{probgrid}{|Y|Y|}
\pcell{11}{Write the LSRL, with the variable names.}{1.7cm}{} &
\pcell{12}{Interpret the slope in context.}{1.7cm}{} \\ \hline
\pcell{13}{Interpret the intercept. Is it meaningful here?}{1.7cm}{} &
\pcell{14}{Find $r$, then interpret $r^2$ in context.}{1.7cm}{} \\ \hline
\end{probgrid}
\\ \hline
\end{guidednotes}

% ── THE NOTES END AT GUIDED PRACTICE ─────────────────────────────────────────
% There is deliberately no "you do" here: ../homework/ is the individual practice,
% started in class in the last 8 minutes while the teacher circulates.

\end{document}
